\documentclass[11pt,a4paper]{article}
\usepackage[utf8]{inputenc}
\usepackage[T1]{fontenc}
\usepackage[english]{babel}
\usepackage{amsmath,amssymb,mathtools}
\usepackage{booktabs}
\usepackage{graphicx}
\usepackage{hyperref}
\usepackage[margin=2.5cm]{geometry}
\usepackage{natbib}
\usepackage{enumitem}
\usepackage{array}
\usepackage{longtable}
\usepackage{xcolor}
\usepackage{microtype}

\hypersetup{
  colorlinks  = true,
  linkcolor   = blue!60!black,
  citecolor   = blue!60!black,
  urlcolor    = blue!60!black,
}

\newcommand{\BAC}{\mathrm{BAC}}
\newcommand{\ie}{\textit{i.e.}}
\newcommand{\eg}{\textit{e.g.}}

\title{Geometry-Survival Evidence Gates for Reproducible and Cost-Aware
       Quantum Kernel Applications Under Noise}

\author{Marcelo Claro Laranjeira\\[4pt]
        \textit{Projeto Fundamentos Qu\^anticos Aplicados \`a Pesquisa em IA}\\[2pt]
        \textit{Protocol frozen: 16 August 2026}\\[2pt]
        SHA-256:\ 89c08f07f2f03a56c144bd4c4b7eddb71813bb1b043bb90920de8e5e70140eaf}

\date{}

\begin{document}
\maketitle

\begin{abstract}
The deployment of quantum kernel methods in machine learning faces a persistent
gap between theoretical promise and empirical validation under realistic noise
conditions. While quantum kernels computed via feature maps on quantum hardware
may capture geometric structures inaccessible to classical kernels, the
conditions under which this potential translates into predictive superiority
remain poorly specified. This study introduces the
\textbf{Geometry-Survival Evidence Gate}, a prospective, sequential decision
framework that integrates paired corrected effect estimation, equivalence
testing, geometric survival analysis across fidelity tiers (ideal statevector,
shot-based estimation, and depolarised noise), and computational cost
accounting before authorising progression to full quantum processing unit (QPU)
classification. The framework was pre-specified via SHA-256 hash and evaluated
on four tabular datasets (synthetic moons, Iris binary, Wine binary, and Breast
Cancer Wisconsin) using nested repeated cross-validation (4 outer folds
$\times$ 3 repetitions) with Nadeau--Bengio correction for dependent folds. The
primary outcome---paired difference in balanced accuracy between a fidelity
quantum kernel SVM and a tuned SVM-RBF---yielded a mean $\Delta\BAC$ of
$-0.217$ (95\% corrected CI $[-0.285;\; -0.148]$; one-sided corrected
$p=0.9999$), indicating statistical inferiority of the quantum kernel under
this protocol. Equivalence testing (TOST, margin $\pm 0.02$) yielded
$p=0.9999$, and the permutation test of signs yielded $p=0.0005$. The
effective kernel rank was 27.9\% of capacity, and kernel--target alignment was
0.110. The validity ladder demonstrated geometric survival across all three
fidelity tiers (BAC stable at 0.625), yet this survival did not translate into
predictive advantage. The multi-dataset exploratory suite confirmed consistent
RBF superiority across all 16 folds (mean $\Delta\BAC = -0.146$). The Evidence
Gate classified this protocol as \textbf{inconclusive for advancement to QPU},
restricting future hypotheses and providing a transparent, reproducible negative
result for quantum kernel evaluation on tabular data.

\medskip
\noindent\textbf{Keywords:} quantum kernel methods; machine learning; nested
cross-validation; evidence gate; negative result; reproducibility; noise
characterization
\end{abstract}

% ── RESUMO ────────────────────────────────────────────────────────────────────
\selectlanguage{brazil}
\begin{abstract}
O emprego de m\'etodos de kernel qu\^antico em aprendizado de m\'aquina enfrenta
uma lacuna persistente entre a promessa te\'orica e a valida\c{c}\~ao
emp\'irica sob condi\c{c}\~oes realistas de ru\'ido. Embora kernels qu\^anticos
computados via feature maps em hardware qu\^antico possam capturar
geometrias inacess\'iveis a kernels cl\'assicos, as condi\c{c}\~oes sob as
quais esse potencial se traduz em superioridade preditiva permanecem mal
especificadas. Este estudo introduz o \textbf{Geometry-Survival Evidence
Gate}---framework prospectivo e sequencial de decis\~ao que integra
estimativa de efeito pareado corrigido, teste de equival\^encia, an\'alise de
sobreviv\^encia geom\'etrica em camadas de fidelidade (statevector ideal,
estimativa por shots e ru\'ido despolarizante) e contabilidade de custo
computacional antes de autorizar o avan\c{c}o para classifica\c{c}\~ao
completa em processador qu\^antico (QPU). O protocolo foi pr\'e-registrado via
hash SHA-256 e avaliado em quatro bases de dados tabulares (moons sint\'etico,
Iris bin\'ario, Wine bin\'ario e Breast Cancer Wisconsin) utilizando
valida\c{c}\~ao cruzada aninhada repetida (4 folds externos $\times$ 3
repeti\c{c}\~oes) com corre\c{c}\~ao de Nadeau--Bengio para folds dependentes.
O desfecho prim\'ario---diferen\c{c}a pareada de acur\'acia balanceada entre
SVM com kernel qu\^antico de fidelidade e SVM-RBF sintonizado---produziu
$\Delta\BAC$ m\'edio de $-0{,}217$ (IC95\% corrigido $[-0{,}285;\;
-0{,}148]$; $p$ unicaudal corrigido $= 0{,}9999$), indicando inferioridade
estat\'istica do kernel qu\^antico neste protocolo. O teste de equival\^encia
(TOST, margem $\pm 0{,}02$) retornou $p = 0{,}9999$ e o teste de
permuta\c{c}\~ao de sinais retornou $p = 0{,}0005$. O posto efetivo do kernel
foi de 27,9\% da capacidade e o alinhamento kernel-alvo foi de 0,110. A escada
de validade demonstrou sobreviv\^encia geom\'etrica nas tr\^es camadas de
fidelidade (BAC est\'avel em 0,625), por\'em essa sobreviv\^encia n\~ao se
converteu em vantagem preditiva. A su\'ite explorat\'oria multibase confirmou
superioridade consistente do RBF em todos os 16 folds ($\Delta\BAC$ m\'edio
$= -0{,}146$). O Evidence Gate classificou este protocolo como
\textbf{inconclusivo para avan\c{c}o a QPU}, restringindo hip\'oteses para
dados tabulares com feature maps de baixa capacidade e demonstrando o Evidence
Gate como mecanismo de controle de qualidade para pesquisa em QML.

\medskip
\noindent\textbf{Palavras-chave:} m\'etodos de kernel qu\^antico; aprendizado
de m\'aquina; valida\c{c}\~ao cruzada aninhada; gate de evid\^encia; resultado
negativo; reprodutibilidade; caracteriza\c{c}\~ao de ru\'ido
\end{abstract}
\selectlanguage{english}

% ══════════════════════════════════════════════════════════════════════════════
\section{Introduction}
\label{sec:intro}

Quantum machine learning (QML) has been presented as one of the most promising
applications of near-term quantum computing, with claims that quantum kernels
could capture geometric structures in high-dimensional spaces inaccessible to
classical approaches \citep{bowles2024better,heyrand2022noisy}. The central
premise is that a quantum feature map---a parametrised sequence of quantum gates
that encodes classical data into quantum states---can produce similarity
matrices (kernels) that, when fed into a classical classifier such as the
Support Vector Machine (SVM), offer predictive advantage over widely used
classical kernels such as the Radial Basis Function (RBF)
\citep{kakavand2026benchmarking}.

However, recent literature has revealed a significant gap between this promise
and the available empirical evidence. \citet{bowles2024better} demonstrated
that many QML benchmarks lack strong baselines and adequate control of
simplistic claims. \citet{kakavand2026benchmarking} conducted the most
comprehensive study to date---with 970 experiments, nested cross-validation,
and IBM hardware validation---and concluded that quantum kernels for tabular
data rarely outperform well-tuned classical baselines.
\citet{heyrand2022noisy} characterised the geometric distortion induced by
noise in quantum kernel machines, showing that the effective rank and spectrum
of the kernel matrix are substantially altered by physical noise channels.
\citet{yin2025experimental} demonstrated superiority on a photonic processor,
but on a platform and task that differs from the standard tabular flow.
\citet{sahin2025quantum} proposed approximation techniques (KTA and
Nystr\"om) to reduce evaluation cost, but without integrating paired
corrected inference into the decision flow.
\citet{aqka2026active} optimised anchor-pair acquisition and shot budgets, but
presupposes that progression to QPU has already been authorised.

In this landscape, a clear methodological gap emerges: no prospective and
sequential rule currently exists that integrates criteria of statistical effect,
equivalence, geometric survival through fidelity tiers, computational cost, and
hardware validation into a transparent go/no-go decision structure. Existing
studies treat these criteria in a fragmented manner---either focusing on
benchmarks without corrected inference, or spectral diagnostics without fidelity
progression, or cost optimisation without a decision gate.

The present article contributes the \textbf{Geometry-Survival Evidence
Gate}---a pre-specified decision rule that answers the following operational
question: \textit{should one proceed from simulation to quantum hardware on a
given classification task?} The Evidence Gate integrates five sequential
criteria: (i) statistical superiority of the quantum kernel over the classical
baseline, assessed by Nadeau--Bengio corrected testing; (ii) practical
equivalence within a pre-specified margin; (iii) survival of accuracy across
three fidelity tiers---exact statevector, 2048-shot estimation, and Aer
depolarised noise simulation; (iv) complete computational cost accounting; and
(v) anchor-pair validation with minimum-target correlation between ideal
fidelity and QPU. The rule is binary (go/no-go) and was frozen via SHA-256
hash before experiment execution.

This study has two declared objectives. The first is to present and formalise
the Evidence Gate as a reusable methodological contribution for the QML
community. The second is to honestly report the negative result obtained: under
the pre-specified protocol, the SVM with fidelity quantum kernel demonstrated
statistical inferiority to the SVM-RBF on tabular data, with $\Delta\BAC =
-0.217$ (95\% corrected CI $[-0.285;\; -0.148]$). This result is
restrictive---it does not settle the question of quantum kernel utility in
general, but constrains future hypotheses and demonstrates the applicability of
the Evidence Gate as a quality control mechanism for QML research.

Section~\ref{sec:theory} presents the theoretical background on quantum
kernels, nested cross-validation, and corrected statistical inference.
Section~\ref{sec:methods} describes the experimental protocol.
Section~\ref{sec:results} reports the results.
Section~\ref{sec:discussion} discusses implications, limitations, and future
work. Section~\ref{sec:conclusions} concludes.

% ══════════════════════════════════════════════════════════════════════════════
\section{Theoretical Background}
\label{sec:theory}

\subsection{Quantum kernels and feature maps}
\label{ssec:kernels}

A quantum kernel is defined by the similarity function between two data points,
computed via the inner overlap of quantum states prepared by a parametrised
circuit \citep{kakavand2026benchmarking}:
\begin{equation}
  K(\mathbf{x},\mathbf{y})
  = \bigl|\langle \phi(\mathbf{x}) | \phi(\mathbf{y})\rangle\bigr|^{2},
  \label{eq:kernel}
\end{equation}
where $|\phi(\mathbf{x})\rangle$ is the quantum state prepared by the feature
map from the classical data vector $\mathbf{x}$. The ZZFeatureMap used in this
study applies $R_z$ rotations with angular data encoding followed by
controlled-Z entangling gates between pairs of qubits, producing interference
that captures pairwise feature correlations \citep{havlicek2019supervised}. For
two qubits with one repetition and linear entanglement, the circuit is:
\begin{equation}
  U_{\phi}(\mathbf{x})
  = \exp\!\Bigl(i \textstyle\sum_{S} \phi_{S}(\mathbf{x})
      \prod_{k\in S} Z_{k}\Bigr)\; H^{\otimes n},
  \label{eq:featuremap}
\end{equation}
where $S$ indexes subsets of qubits with entanglement and $\phi_S$ are data
encoding functions. The kernel matrix $K$ is evaluated by computing
$|\langle \phi(\mathbf{x}_i)|\phi(\mathbf{x}_j)\rangle|^{2}$ for all pairs
$(i,j)$ in the dataset, producing a positive semi-definite matrix that feeds
the SVM \citep{havlicek2019supervised,kakavand2026benchmarking}.

Exact evaluation of the kernel matrix requires inner overlap computation via
statevector (perfect quantum simulation). On real hardware, the estimate is
obtained by sampling (shots): each data pair is prepared as a quantum state,
measured repeatedly, and fidelity is estimated by outcome counting
\citep{heyrand2022noisy}. The variance of this estimate scales as
$O(1/N_{\text{shots}})$, introducing statistical noise that adds to
hardware physical noise.

\subsection{Noise and geometric degradation}
\label{ssec:noise}

On real quantum hardware, three error sources affect the kernel matrix: (i)
gate errors, which distort the prepared state; (ii) readout errors, which
alter measured outcomes; and (iii) decoherence, which destroys quantum
correlations \citep{heyrand2022noisy}. Modelable noise channels---such as the
depolarising channel with parameter $p$---transform the ideal state
$\rho_0$ into:
\begin{equation}
  \mathcal{E}(\rho) = (1-p)\,\rho_0 + \frac{p}{2^{n}}\,I,
  \label{eq:depolarise}
\end{equation}
where $n$ is the number of qubits and $I$ is the identity matrix. This
distortion alters the kernel matrix spectrum, reducing kernel--target alignment
and modifying the effective rank \citep{heyrand2022noisy,sahin2025quantum}.

The concept of \textbf{geometric survival}---the maintenance of predictive
accuracy across successive fidelity tiers
(statevector$\;\to\;$shots$\;\to\;$noise)---is central to this study.
Survival is not guaranteed: shot statistical fluctuations can alter
classification decisions, and physical noise can degrade the kernel structure
\citep{heyrand2022noisy}. The validity ladder quantifies this survival
empirically.

\subsection{Nested cross-validation and corrected inference}
\label{ssec:cv}

Cross-validation (CV) is the standard for generalisation estimation in machine
learning \citep{refsgaard1999proper}. However, standard CV violates the
independence assumption between folds, inflating statistical power and
producing underestimated confidence intervals \citep{nadeau2003inference}.
Nested CV partially addresses this problem by separating hyperparameter
selection (inner folds) from generalisation estimation (outer folds)
\citep{refsgaard1999proper}.

\citet{nadeau2003inference} proposed a correction for the standard error in
dependent CV estimates:
\begin{equation}
  \mathrm{SE}_{\mathrm{corr}}
  = \mathrm{SE} \times
    \sqrt{\frac{1}{n_{\text{test}}} + \frac{n_{\text{test}}}{n_{\text{train}}}},
  \label{eq:nadeau}
\end{equation}
where $n_{\text{test}}$ and $n_{\text{train}}$ are the test and training set
sizes, respectively. This correction is conservative: it produces wider
confidence intervals that adequately reflect the dependence between folds.

The primary one-sided test---$H_1$: paired mean $\Delta\BAC > 0$---is
evaluated with the Nadeau--Bengio corrected $t$-statistic. Sensitivity analyses
include: (i) exact sign permutation, testing $H_0$:
$\mathrm{median}(\Delta)=0$ without normality assumption; (ii) TOST
(Two One-Sided Tests) equivalence with pre-specified margin $\pm 0.02$ in BAC;
and (iii) Holm adjustment for secondary outcomes.

\subsection{Evidence gates in reproducible science}
\label{ssec:evidencegates}

The concept of an evidence gate has roots in the clinical trials literature and
open science practices, where go/no-go criteria are pre-specified before data
collection \citep{scheel2022role}. The Evidence Gate proposed in this study
adapts this concept to QML, integrating statistical, geometric, and economic
criteria into a formal decision sequence.

The distinction between the Evidence Gate and existing approaches is important.
\citet{kakavand2026benchmarking} conducted rigorous benchmarks, but without a
formal decision rule for when to advance to hardware.
\citet{aqka2026active} optimised resource allocation, but presupposes prior
advancement. \citet{bowles2024better} cautioned against simplistic claims, but
did not provide an operational decision structure. The Evidence Gate fills this
gap by formalising the question: \textit{based on simulated evidence, should
one reserve QPU time for this task?}

% ══════════════════════════════════════════════════════════════════════════════
\section{Methods}
\label{sec:methods}

\subsection{Experimental design}
\label{ssec:design}

The study was designed as a comparative computational investigation,
pre-specified and frozen via SHA-256 hash before execution. The primary outcome
was the paired difference in balanced accuracy ($\Delta\BAC = \BAC_{\mathrm{QML}} -
\BAC_{\mathrm{RBF}}$). The pre-specified classification of the result was:
superiority, inferiority, practical equivalence, or inconclusive.

Four tabular datasets were used: (i) \texttt{make\_moons} (synthetic, 200
samples, 2 features, Gaussian noise $\sigma=0.1$); (ii) Iris binary (100
samples, 4 features, classes 0/1); (iii) Wine binary (130 samples, 13
features, classes 0/1); (iv) Breast Cancer Wisconsin (180 samples, 30
features, classes malignant/benign). All datasets were binarised and subjected
to identical preprocessing: missing-value imputation (median), standardisation
($z$-score), and dimensionality reduction by PCA to 2 components. All
preprocessing was fit exclusively on the training set in each fold to prevent
data leakage.

\subsection{Repeated nested cross-validation}
\label{ssec:ncv}

The validation design used repeated nested CV with: (i) 4 stratified outer
folds (generalisation estimation); (ii) 3 outer repetitions with different
seeds; (iii) 3 stratified inner folds (hyperparameter selection). The total
was 12 outer evaluations per dataset.

In each outer fold, the imputation--standardisation--PCA--angular scaling
transformation was fit exclusively on the training data. Internal selection
used a hyperparameter grid: $C \in \{0.1,\; 1.0,\; 10.0\}$ for both SVMs and
$\gamma \in \{\text{scale},\; \text{auto}\}$ for the RBF. The winning internal
classifier was selected by mean balanced accuracy on inner folds and then
re-evaluated on the outer fold.

\subsection{Quantum kernel}
\label{ssec:qkernel}

The quantum feature map used ZZFeatureMap with 2 qubits, 1 repetition, and
linear entanglement. The kernel matrix was evaluated across three fidelity tiers
(validity ladder):
\begin{enumerate}[label=\arabic*.]
  \item \textbf{Exact statevector:} direct inner overlap computation without
        sampling---represents the ideal noise-free limit.
  \item \textbf{2048 shots:} sampling-based estimation with 2048 independent
        measurements per pair---introduces statistical fluctuation.
  \item \textbf{Aer with noise:} simulation of the Qiskit Aer depolarising
        channel with parameters $p_1 = 0.001$ (single-qubit gate error),
        $p_2 = 0.01$ (two-qubit gate error) and $p_{\text{readout}} = 0.02$
        (readout error)---represents typical NISQ hardware conditions.
\end{enumerate}

\subsection{Evidence Gate---decision rule}
\label{ssec:gate}

The Evidence Gate is a sequential decision rule with five pre-specified
criteria. The protocol advances through stages only if \textbf{all} criteria
of the current stage are satisfied:

\medskip
\noindent\textbf{Stage 1 --- Primary inference:}
\begin{itemize}
  \item $\Delta\BAC$ corrected mean $> 0$ (statistical superiority)
  \item 95\% bilateral corrected CI (Nadeau--Bengio) above zero
  \item Classification: inferiority, equivalence, or inconclusive
        $\to$ \textbf{stop}
\end{itemize}

\noindent\textbf{Stage 2 --- Practical equivalence:}
\begin{itemize}
  \item TOST with margin $\pm 0.02$ in BAC yields $p < 0.05$
  \item If not equivalent $\to$ continue only if Stage 3 is promising
\end{itemize}

\noindent\textbf{Stage 3 --- Geometric survival:}
\begin{itemize}
  \item BAC at exact statevector $\geq$ classical baseline
  \item BAC with 2048 shots $\geq 95\%$ of statevector BAC
  \item BAC with Aer noise $\geq 90\%$ of statevector BAC
  \item Drop $>10\%$ across any transition $\to$ \textbf{stop}
\end{itemize}

\noindent\textbf{Stage 4 --- Computational cost:}
\begin{itemize}
  \item Kernel time recorded and comparable to baseline
  \item Estimated logical evaluations reported
  \item Infeasible cost $\to$ \textbf{stop}
\end{itemize}

\noindent\textbf{Stage 5 --- Anchor pairs and QPU advancement:}
\begin{itemize}
  \item Pearson/Spearman correlation ideal--QPU $\geq 0.90$ on anchor pairs
  \item MAE $\leq 0.10$ between ideal fidelity and QPU on anchor pairs
  \item Budgetary cost viability
  \item All satisfied $\to$ \textbf{go} for full QPU classifier
\end{itemize}

\subsection{Statistical inference}
\label{ssec:statistics}

The primary test used Nadeau--Bengio correction for fold dependence
\citep{nadeau2003inference}, with $\alpha = 0.05$ and one-sided test ($H_1$:
$\mu_{\Delta\BAC} > 0$). Pre-specified sensitivity analyses: (i) exact sign
permutation with 1024 permutations; (ii) TOST equivalence with margin
$\pm 0.02$; (iii) paired Cohen's $d_z$; (iv) 95\% bootstrap interval (5000
resamples); (v) Holm adjustment for secondary outcomes.

\subsection{Kernel diagnostics}
\label{ssec:diagnostics}

Three diagnostic metrics were recorded at each evaluation: (i) kernel--target
alignment \citep{bowles2024better}; (ii) effective rank, defined as
$\exp(-\sum_i p_i \ln p_i)$ where $p_i = \lambda_i / \sum_j \lambda_j$
\citep{heyrand2022noisy}; (iii) relative Frobenius error between the evaluated
kernel matrix and the statevector reference.

\subsection{Claim control}
\label{ssec:claims}

Claim control rules were pre-specified \citep{bowles2024better,kakavand2026benchmarking}:
\begin{itemize}
  \item Do not claim ``first study'' without an updated systematic review.
  \item Do not claim ``quantum advantage'' from simulator accuracy.
  \item Distinguish predictive superiority, practical equivalence, geometric
        survival, and computational advantage.
  \item Report null and negative results without softening.
  \item Identify as exploratory any analysis not frozen in the SHA-256 protocol.
\end{itemize}

% ══════════════════════════════════════════════════════════════════════════════
\section{Results}
\label{sec:results}

\subsection{Primary result---statistical inferiority}
\label{ssec:primary}

In repeated nested cross-validation (12 outer evaluations,
\texttt{make\_moons} dataset), the mean primary outcome was $\Delta\BAC =
-0.217$, with Nadeau--Bengio corrected 95\% CI $[-0.285;\; -0.148]$ and
corrected one-sided $p=0.9999$ (Table~\ref{tab:inference}). The confidence
interval is entirely negative, indicating that the SVM-RBF outperforms the
fidelity quantum kernel in balanced accuracy within this protocol.

\begin{table}[htbp]
\centering
\caption{Statistical inference results for the primary and secondary outcomes.}
\label{tab:inference}
\smallskip
\begin{tabular}{@{}llccc@{}}
\toprule
Outcome & Mean & 95\% CI Lower & 95\% CI Upper & Corrected $p$ \\
\midrule
$\Delta$ balanced accuracy (primary) & $-0.217$ & $-0.285$ & $-0.148$ & $0.9999$ \\
$\Delta$ accuracy                      & $-0.217$ & $-0.285$ & $-0.148$ & $2.38\times10^{-5}$ \\
$\Delta F1$                            & $-0.200$ & $-0.296$ & $-0.104$ & $8.05\times10^{-4}$ \\
\bottomrule
\end{tabular}
\end{table}

The paired Cohen's $d_z = -4.495$, indicating a large-magnitude difference
between models. This value reflects the consistency of the pattern: across all
12 outer evaluations, the SVM-RBF outperformed the quantum kernel.

\subsection{Sensitivity analyses}
\label{ssec:sensitivity}

Exact sign permutation (1024 permutations) yielded $p=0.0005$, confirming
statistical inferiority. TOST equivalence with margin $\pm 0.02$ yielded
$p=0.9999$, indicating the models are \emph{not} equivalent.

\begin{table}[htbp]
\centering
\caption{Pre-specified sensitivity analyses.}
\label{tab:sensitivity}
\smallskip
\begin{tabular}{@{}llll@{}}
\toprule
Analysis & Statistic & $p$ / CI & Interpretation \\
\midrule
Nadeau--Bengio (one-sided $H_1$: $\Delta>0$) & Corrected $t$ & $p=0.9999$ & No superiority evidence \\
Sign permutation (bilateral)  & Exact $p$  & $p=0.0005$ & Significant inferiority \\
TOST equivalence ($\pm 0.02$) & One-sided $p$ & $p=0.9999$ & Not equivalent \\
Cohen's $d_z$ & $d_z=-4.495$ & --- & Large negative effect \\
95\% bootstrap $\Delta\BAC$ (5000) & Median & $[-0.563;\; 0.063]$ & Includes zero (uncorrected) \\
Fraction positive $\Delta$ resamples & --- & 4.2\% & Rarely does quantum kernel win \\
\bottomrule
\end{tabular}
\end{table}

\subsection{Robust validation---outer repetitions}
\label{ssec:repetitions}

Table~\ref{tab:ncv} presents the 12 outer evaluations.

\begin{table}[htbp]
\centering
\caption{Repeated nested cross-validation results (12 outer evaluations).}
\label{tab:ncv}
\smallskip
\footnotesize
\begin{tabular}{@{}ccccccccc@{}}
\toprule
Run & Rep. & Fold & Seed & $\BAC$ RBF & $\BAC$ QML & $\Delta\BAC$ & $C_{\text{RBF}}$ & $C_{\text{QML}}$ & Kernel (s) \\
\midrule
 1 & 1 & 1 & 1051  & 1.000 & 0.767 & $-0.233$ & 10.0 &  1.0 & 0.89 \\
 2 & 1 & 2 & 2060  & 0.933 & 0.733 & $-0.200$ & 10.0 &  0.1 & 0.91 \\
 3 & 1 & 3 & 3069  & 0.967 & 0.767 & $-0.200$ & 10.0 &  0.1 & 0.92 \\
 4 & 1 & 4 & 4078  & 0.933 & 0.800 & $-0.133$ & 10.0 &  0.1 & 0.91 \\
 5 & 2 & 1 & 5087  & 0.967 & 0.767 & $-0.200$ & 10.0 &  1.0 & 0.96 \\
 6 & 2 & 2 & 6096  & 0.967 & 0.667 & $-0.300$ & 10.0 & 10.0 & 0.95 \\
 7 & 2 & 3 & 7105  & 0.900 & 0.667 & $-0.233$ & 10.0 &  0.1 & 1.12 \\
 8 & 2 & 4 & 8114  & 0.967 & 0.733 & $-0.233$ &  1.0 &  1.0 & 1.29 \\
 9 & 3 & 1 & 9123  & 1.000 & 0.700 & $-0.300$ & 10.0 & 10.0 & 1.44 \\
10 & 3 & 2 & 10132 & 0.833 & 0.667 & $-0.167$ &  1.0 & 10.0 & 0.94 \\
11 & 3 & 3 & 11141 & 1.000 & 0.800 & $-0.200$ & 10.0 & 10.0 & 0.89 \\
12 & 3 & 4 & 12150 & 0.933 & 0.733 & $-0.200$ & 10.0 &  1.0 & 0.95 \\
\bottomrule
\end{tabular}
\end{table}

The SVM-RBF won in 12/12 evaluations (100\%), with mean BAC of 0.955 (SD
0.050). The quantum kernel showed mean BAC of 0.738 (SD 0.047). The mean
$\Delta\BAC$ was $-0.217$ (SD 0.048).

\subsection{Validity ladder---geometric survival}
\label{ssec:ladder}

\begin{table}[htbp]
\centering
\caption{Validity ladder---statevector $\to$ shots $\to$ Aer noise.}
\label{tab:ladder}
\smallskip
\footnotesize
\begin{tabular}{@{}lccccccc@{}}
\toprule
Level & Accuracy & $F1$ & Frobenius & Alignment & Eff.\ Rank & Time (s) & Errors \\
\midrule
0 \textperiodcentered\ Exact SV  & 0.625 & 0.500 & 0.000 & 0.109 & 8.63  & 0.055  & --- \\
1 \textperiodcentered\ 2048 shots & 0.625 & 0.500 & 0.011 & 0.110 & 8.93  & 23.271 & --- \\
2 \textperiodcentered\ Aer noise   & 0.625 & 0.500 & 0.055 & 0.110 & 10.34 & 8.595  & $0.001/0.01/0.02$ \\
\bottomrule
\end{tabular}
\end{table}

Geometric survival was \textbf{complete}: BAC remained stable at 0.625 across
all tiers. Relative Frobenius error grew progressively ($0.000 \to 0.011
\to 0.055$), indicating increasing kernel distortion, but this did not alter
SVM classification decisions.

\subsection{Computational cost}
\label{ssec:cost}

Kernel time ranged from 0.89\,s to 1.44\,s (mean 1.01\,s). The 2048-shot
evaluation required 23.27\,s ($423\times$ statevector), and Aer noise required
8.59\,s ($156\times$ statevector). Estimated logical evaluations: $1{,}008
\times 2048 = 2{,}064{,}384$ shots.

\subsection{Model comparison}
\label{ssec:models}

\begin{table}[htbp]
\centering
\caption{Model comparison---\texttt{make\_moons} (single split).}
\label{tab:models}
\smallskip
\begin{tabular}{@{}lccccc@{}}
\toprule
Model & Accuracy & $\BAC$ & $F1$ & Kernel (s) & 95\% Bootstrap CI \\
\midrule
SVM-RBF                & 0.875 & 0.875 & 0.889 & ---     & $[0.688;\; 1.000]$ \\
Logistic Regression    & 0.813 & 0.813 & 0.842 & ---     & $[0.625;\; 1.000]$ \\
SVM + Quantum Kernel   & 0.625 & 0.625 & 0.500 & 23.27   & $[0.375;\; 0.875]$ \\
\bottomrule
\end{tabular}
\end{table}

\subsection{Multi-dataset suite---exploratory analysis}
\label{ssec:multidataset}

\begin{table}[htbp]
\centering
\caption{Multi-dataset application suite (exploratory, not frozen).}
\label{tab:multi}
\smallskip
\footnotesize
\begin{tabular}{@{}lccccccc@{}}
\toprule
Dataset & Fold & $\Delta\BAC$ & Geom.\ Surv. & Frobenius & Alignment & Rel.\ Rank & Kernel (s) \\
\midrule
make\_moons   & 1 & $-0.300$ & 0.978 & 0.022 & 0.093 & 0.117 & 0.86 \\
make\_moons   & 2 & $-0.200$ & 0.980 & 0.020 & 0.105 & 0.115 & 0.89 \\
make\_moons   & 3 & $-0.133$ & 0.980 & 0.020 & 0.122 & 0.112 & 0.89 \\
make\_moons   & 4 & $-0.167$ & 0.976 & 0.024 & 0.093 & 0.111 & 0.90 \\
iris\_binaria & 1 & $-0.119$ & 0.979 & 0.021 & 0.130 & 0.115 & 0.68 \\
iris\_binaria & 2 & $-0.237$ & 0.936 & 0.064 & 0.147 & 0.127 & 0.63 \\
iris\_binaria & 3 & $-0.080$ & 0.981 & 0.019 & 0.134 & 0.116 & 0.64 \\
iris\_binaria & 4 & $-0.237$ & 0.983 & 0.017 & 0.166 & 0.116 & 0.66 \\
wine\_binario & 1 & $-0.072$ & 0.978 & 0.022 & 0.167 & 0.099 & 1.03 \\
wine\_binario & 2 & $-0.028$ & 0.978 & 0.022 & 0.261 & 0.095 & 1.40 \\
wine\_binario & 3 & $-0.059$ & 0.984 & 0.016 & 0.224 & 0.096 & 1.77 \\
wine\_binario & 4 & $-0.127$ & 0.975 & 0.025 & 0.128 & 0.109 & 1.05 \\
breast\_cancer & 1 & $-0.133$ & 0.984 & 0.016 & 0.189 & 0.074 & 1.55 \\
breast\_cancer & 2 & $-0.020$ & 0.985 & 0.015 & 0.200 & 0.073 & 1.53 \\
breast\_cancer & 3 & $-0.087$ & 0.985 & 0.015 & 0.199 & 0.074 & 1.52 \\
breast\_cancer & 4 & $-0.053$ & 0.984 & 0.016 & 0.195 & 0.078 & 1.49 \\
\bottomrule
\end{tabular}
\end{table}

Across all 16 folds, $\Delta\BAC$ was negative. Geometric survival was high
($>0.93$), but did not translate into predictive advantage.

\subsection{Mechanistic analysis}
\label{ssec:correlation}

\begin{table}[htbp]
\centering
\caption{Exploratory Spearman correlations between kernel metrics and
         $\Delta\BAC$.}
\label{tab:correlation}
\smallskip
\begin{tabular}{@{}llll@{}}
\toprule
Metric & $\rho$ (Spearman) & $p$ & Interpretation \\
\midrule
Geometric survival $\to$ $\Delta\BAC$     & 0.415 & 0.110 & Moderate, non-significant \\
Kernel--target alignment $\to$ $\Delta\BAC$ & 0.739 & 0.001 & Strong, significant \\
\bottomrule
\end{tabular}
\end{table}

\subsection{Kernel diagnostics}
\label{ssec:kdiag}

\begin{table}[htbp]
\centering
\caption{Quantum kernel diagnostics.}
\label{tab:kdiag}
\smallskip
\begin{tabular}{@{}ll@{}}
\toprule
Metric & Value \\
\midrule
Maximum asymmetry (eigenvalues)   & $2.22\times10^{-16}$ \\
Maximum diagonal deviation       & 0.006 \\
Smallest eigenvalue              & $-5.97\times10^{-16}$ \\
Kernel--target alignment         & 0.110 \\
Diagonal tolerance               & 0.022 \\
\bottomrule
\end{tabular}
\end{table}

\subsection{Evidence Gate result}
\label{ssec:egate}

\begin{table}[htbp]
\centering
\caption{Evidence Gate sequential evaluation.}
\label{tab:egate}
\smallskip
\begin{tabular}{@{}p{3.2cm}cp{5.5cm}@{}}
\toprule
Criterion & Status & Result \\
\midrule
\textbf{Stage 1}: $\Delta\BAC>0$, CI$_{95\%}>0$
  & \textbf{FAIL} & $\Delta\BAC=-0.217$; CI $[-0.285;\; -0.148]$ \\
\textbf{Stage 2}: Equivalence (TOST, $\pm 0.02$)
  & \textbf{FAIL} & $p=0.9999$ (not equivalent) \\
\textbf{Stage 3}: Geometric survival
  & \textbf{PASS} & BAC stable 0.625 across 3 tiers \\
\textbf{Stage 4}: Computational cost
  & \textbf{PASS} & Kernel $\sim 1$\,s; cost recorded \\
\textbf{Stage 5}: Anchor-pair QPU
  & \textbf{NOT ASSESSED} & Conditional on Stage 1 \\
\bottomrule
\end{tabular}
\end{table}

\noindent\textbf{Evidence Gate classification: NO-GO---inconclusive for QPU
advancement.}

% ══════════════════════════════════════════════════════════════════════════════
\section{Discussion}
\label{sec:discussion}

\subsection{Statistical inferiority}
\label{ssec:disc_inferiority}

The primary outcome---statistical inferiority of the fidelity quantum kernel
over the SVM-RBF ($\Delta\BAC = -0.217$)---is consistent with recent
literature \citep{bowles2024better,kakavand2026benchmarking}. The RBF
baseline was tuned by nested validation, representing a strong classical
competitor. The quantum kernel showed consistently low kernel--target alignment
(0.093--0.261). This geometric inadequacy---not noise instability---appears
to be the dominant factor in inferiority.

The spectral diagnostics (Table~\ref{tab:kdiag}) reveal that the 2-qubit
feature map may be structurally insufficient to capture class separability in
tabular data with 2--30 features. The low relative effective rank (7--12\%)
indicates underutilisation of circuit capacity.

\subsection{Geometric survival without predictive advantage}
\label{ssec:disc_survival}

A conceptual contribution is the empirical demonstration that
\textbf{geometric survival $\neq$ predictive advantage}. The kernel survived
perfectly across statevector$\;\to\;$shots$\;\to\;$Aer, but its performance
was consistently inferior to the RBF. Survival indicates noise robustness;
inferiority indicates geometric inadequacy.

\subsection{Correlation between alignment and performance}
\label{ssec:disc_alignment}

The strong correlation between kernel--target alignment and $\Delta\BAC$
($\rho = 0.739$, $p = 0.001$) is consistent with kernel theory. However, even
the highest observed alignment (0.261) was insufficient to produce positive
$\Delta\BAC$, suggesting an alignment threshold below which the quantum kernel
offers no advantage.

\subsection{Positioning relative to the literature}
\label{ssec:disc_positioning}

The present study distinguishes itself in three aspects: (i) the Evidence Gate
as a methodological contribution; (ii) transparent negative result
\citep{scheel2022role}; (iii) cost integration as a decision factor.

\subsection{Limitations}
\label{ssec:limitations}

\begin{enumerate}[label=(\roman*)]
  \item Small scale: 200 samples, 90/10 splits.
  \item Limited feature map: 2 qubits, 1 repetition.
  \item Simulation only (Aer), not real QPU.
  \item Tabular data only; may differ for molecular/materials data.
  \item PCA to 2 components may discard discriminative information.
  \item Residual fold dependence despite correction.
  \item Mechanistic analyses are exploratory (not frozen in SHA-256).
\end{enumerate}

\subsection{Future work}
\label{ssec:future}

\begin{enumerate}[label=(\roman*)]
  \item Feature map expansion to identify alignment threshold for advancement.
  \item QPU validation with anchor pairs (Criteria 3--5).
  \item Cost reduction via approximation (Nystr\"om, KTA).
  \item Multi-class and regression extensions.
  \item Quantum-structure data (molecular, materials).
  \item Systematic review with retrospective Evidence Gate application.
\end{enumerate}

\subsection{Implications for QML practice}
\label{ssec:implications}

\begin{enumerate}[label=(\roman*)]
  \item A strong baseline is essential; weak baselines produce superficial
        superiority claims.
  \item Kernel diagnostics (alignment, effective rank, Frobenius error) must
        be routinely reported.
  \item Negative results are valuable and constrain future hypotheses.
  \item Evidence Gates should be adopted as minimum standard for QML
        evaluation.
\end{enumerate}

% ══════════════════════════════════════════════════════════════════════════════
\section{Conclusions}
\label{sec:conclusions}

This study presented the \textbf{Geometry-Survival Evidence Gate}---a
prospective, sequential go/no-go decision rule for quantum kernel application
on tabular data---and honestly reported the negative result obtained under a
pre-specified protocol.

The SVM with fidelity quantum kernel demonstrated \textbf{statistical
inferiority} to the SVM-RBF in repeated nested cross-validation ($\Delta\BAC =
-0.217$; 95\% corrected CI $[-0.285;\; -0.148]$; one-sided $p = 0.9999$).
The kernel survived geometrically across three fidelity tiers (statevector
$\to$ shots $\to$ Aer noise), but this survival did not translate into
predictive advantage. The correlation between kernel--target alignment and
$\Delta\BAC$ ($\rho = 0.739$) suggests that feature map geometry is the
limiting factor, not noise robustness.

The Evidence Gate classified this protocol as \textbf{NO-GO} for QPU
advancement, with failure at Criteria~1 (superiority) and~2 (equivalence). The
result is restrictive: it does not settle quantum kernel utility in general, but
constrains hypotheses for tabular data with low-capacity feature maps and
demonstrates the Evidence Gate as a quality control mechanism for QML research.

It is recommended that the community adopt prospective evidence gates as
minimum standard for quantum kernel evaluation, valuing transparent negative
results and promoting reproducibility in an area prone to premature claims.

% ══════════════════════════════════════════════════════════════════════════════
\bibliographystyle{apalike}
\bibliography{ARTIGO_QUANTUM_EVIDENCE_GATE}

% ══════════════════════════════════════════════════════════════════════════════
\appendix
\section{Experiment Configuration}
\label{app:config}

\begin{table}[htbp]
\centering
\begin{tabular}{@{}ll@{}}
\toprule
Parameter & Value \\
\midrule
Researcher & Marcelo Claro Laranjeira \\
Seed & 42 \\
Shots & 2048 \\
Feature map & ZZFeatureMap \\
Feature map reps & 1 \\
Entanglement & Linear \\
$C$ (quantum SVM) & 1.0 \\
Kernel time (full split) & 23.27\,s \\
Kernel--target alignment & 0.110 \\
Diagonal tolerance & 0.022 \\
Protocol SHA-256 & \texttt{89c08f07\ldots} \\
Python & 3.12.13 \\
NumPy & 2.0.2 \\
Pandas & 2.2.2 \\
Scikit-learn & 1.6.1 \\
Qiskit & 2.3.1 \\
Qiskit Aer & 0.17.2 \\
Qiskit ML & 0.9.0 \\
\bottomrule
\end{tabular}
\end{table}

\section{Pre-registered Protocol}
\label{app:protocol}

The full protocol was frozen via SHA-256 and is available in
\texttt{data/processed/protocolo\_sha256.txt}. The mandatory execution order
was: OSF $\to$ Nested CV $\to$ Application suite $\to$ Aer noise $\to$
Anchor pairs QPU $\to$ Full QPU classifier. The present execution completed
items 1--3 and stopped at Evidence Gate Criterion~1 (statistical inferiority).

\end{document}
